\quad O \textit{Quick Sort} é um algoritmo de ordenação baseado na
estratégia de divisão e conquista. Sua ideia principal é escolher um
elemento do vetor, chamado de pivô, e reorganizar os demais elementos de
modo que os valores menores que o pivô fiquem de um lado e os valores
maiores fiquem do outro.

\quad Após esse particionamento, o algoritmo é aplicado recursivamente às
duas partes formadas, até que cada partição possua no máximo um elemento.
Diferentemente do \textit{Merge Sort}, não é necessário um vetor auxiliar
para intercalar os resultados, ja que a ordenação ocorre no vetor.

\quad Foram implementadas duas versões do algoritmo, que
diferem apenas na função de particionamento utilizada, mais
especificamente na escolha do pivô. Na versão 1 (\texttt{particao\_v1}),
o pivô é o primeiro elemento da partição, obtido diretamente em
\texttt{vetor[esq]}, sem custo adicional. Na versão 2
(\texttt{particao\_v2}), o pivô é a média aritmética dos elementos da
partição, calculada a cada chamada ao percorrer todos os elementos entre
\texttt{esq} e \texttt{dir}, somá-los e dividir o resultado pela
quantidade de elementos. Como a média nem sempre corresponde a um
elemento do vetor, nessa versão o particionamento é feito por comparação
com o valor da média, e não com uma posição específica. Fora isso, as duas funções são idênticas.

\newpage

\begin{lstlisting}[caption={Implementação do Quick Sort}]
void quickSort(int* vetor, int tamanho, int versao) {

    auto ordena = [&](auto&& self, int esq, int dir) -> void {
        while (esq < dir) {
            int i, j;

            if (versao == 1)
                particao_v1(vetor, esq, dir, i, j);
            else if (versao == 2)
                particao_v2(vetor, esq, dir, i, j);

            int tamEsq = j - esq + 1;
            int tamDir = dir - i + 1;

            if (tamEsq < tamDir) {
                if (esq < j)
                    self(self, esq, j);

                esq = i;
            } else {
                if (i < dir)
                    self(self, i, dir);

                dir = j;
            }
        }
    };

    if (tamanho > 1 && (versao == 1 || versao == 2))
        ordena(ordena, 0, tamanho - 1);
}
\end{lstlisting}

\quad Nas duas funções de particionamento, os índices \texttt{i} e
\texttt{j} começam nas extremidades da partição e se aproximam até se
cruzarem. O índice \texttt{i} avança enquanto encontra elementos menores
que o pivô, e o índice \texttt{j} volta enquanto encontra elementos
maiores que o pivô. Quando ambos param, e \texttt{i} ainda não
ultrapassou \texttt{j}, os elementos são trocados de posição e os dois
índices continuam.

\quad No final do particionamento, o índice \texttt{j} indica o fim da
partição da esquerda e o índice \texttt{i} indica o início da partição da
direita, assim, os elementos à esquerda são menores ou iguais
ao pivô e os elementos à direita são maiores ou iguais a ele. A
diferença entre as versões está, portanto, no custo e na qualidade da
escolha do pivô: a versão 1 não exige nenhuma operação extra, enquanto a
versão 2 acrescenta um custo linear a cada particionamento, em troca de
uma divisão potencialmente mais equilibrada.

\quad A função \texttt{quickSort} recebe o vetor, seu tamanho e a versão
desejada. Uma função lambda recursiva (\texttt{ordena})
recebe os limites \texttt{esq} e \texttt{dir} da partição atual, chama a
função de particionamento correspondente à versão escolhida e calcula o
tamanho das duas partições geradas, \texttt{tamEsq} e \texttt{tamDir}.

\quad Para limitar a profundidade da recursão, a chamada recursiva é feita
apenas sobre a menor das duas partições, enquanto a maior é tratada pelo
próprio laço \texttt{while}, ajustando os limites \texttt{esq} ou
\texttt{dir}. O processo termina quando a partição possui no máximo um
elemento.

\quad A escolha do pivô influencia diretamente o balanceamento das
partições. Na versão 1, vetores já ordenados de forma crescente ou
decrescente tendem a gerar partições muito desiguais, aproximando o
algoritmo do pior caso. Na versão 2, a média aritmética tende a dividir
melhor os valores da partição, porém exige a soma dos elementos a cada
chamada, o que acrescenta um custo adicional. Esses efeitos serão
observados nos experimentos apresentados neste trabalho.